\documentclass[11pt,a4paper]{article}

\usepackage[utf8]{inputenc}
\usepackage[T1]{fontenc}
\usepackage{lmodern}
\usepackage[french]{babel}
\usepackage[margin=2.2cm]{geometry}
\usepackage{tabularx}
\usepackage{amssymb}
\usepackage{longtable}
\usepackage{booktabs}
\usepackage{array}
\usepackage{colortbl}
\usepackage{xcolor}
\usepackage{graphicx}
\usepackage{enumitem}
\usepackage{titlesec}
\usepackage[colorlinks=true,linkcolor=blue!60!black,urlcolor=blue!60!black]{hyperref}

\hypersetup{
  pdftitle={PEKEGNO Management System -- Plan d'implementation et decoupage des taches},
  pdfauthor={Equipe developpement Pekegno},
  pdfsubject={Plan de developpement, analyse des ecarts et planning}
}

\setlength{\parskip}{0.35em}

% Couleurs
\definecolor{pekblue}{RGB}{54,65,245}
\definecolor{peklight}{RGB}{228,231,255}
\definecolor{pekgreen}{RGB}{21,150,110}
\definecolor{pekamber}{RGB}{247,144,9}
\definecolor{pekgray}{RGB}{120,130,145}

\newcolumntype{W}{>{\centering\arraybackslash}p{0.40cm}}
\newcolumntype{Y}{>{\raggedright\arraybackslash}X}

\newcommand{\be}{\textsc{Be}}
\newcommand{\fe}{\textsc{Fe}}
\newcommand{\qa}{\textsc{Qa}}

\begin{document}

%=====================================================
% Page de titre
%=====================================================
\begin{titlepage}
  \centering
  \vspace*{2.5cm}
  {\Huge\bfseries PEKEGNO\par}
  \vspace{0.3cm}
  {\LARGE MANAGEMENT SYSTEM\par}
  \vspace{1.8cm}
  {\Large\bfseries Plan d'implémentation\\[0.3cm] Analyse des écarts \& découpage des tâches\par}
  \vspace{1.5cm}
  \begin{tabular}{rl}
    \textbf{Base de travail} & Cahier des charges fonctionnel (Août 2026) \\[2mm]
    \textbf{Périmètre} & Backend Laravel + Frontend React existants \\[2mm]
    \textbf{Équipe} & Blackstar \& Le\_million \\[2mm]
    \textbf{Durée prévisionnelle} & 20 semaines (10 sprints de 2 semaines) \\
  \end{tabular}
  \vfill
\end{titlepage}

\tableofcontents
\newpage

%=====================================================
\section{Objet du document}
%=====================================================

Ce document fait le lien entre le \emph{cahier des charges fonctionnel} PEKEGNO et l'application réellement développée à ce jour (backend Laravel 13 / PostgreSQL, frontend React 19 / Vite / Tailwind v4). Il :

\begin{enumerate}[nosep]
  \item dresse l'état des lieux précis du code existant ;
  \item identifie les écarts entre l'existant et les exigences du cahier des charges ;
  \item corrige deux interprétations erronées validées avec la direction : le modèle \textbf{Agence $\rightarrow$ Départements typés} et la navigation par \textbf{barre de contexte permanente} ;
  \item découpe le travail restant en tâches chiffrées en durée, réparties entre deux développeurs, \textbf{Blackstar} et \textbf{Le\_million}, qui permutent entre backend et frontend à chaque sprint.
\end{enumerate}

%=====================================================
\section{État des lieux du projet}
%=====================================================

\subsection{Backend (Laravel 13, PostgreSQL)}

Environ 70 migrations, 40 modèles Eloquent, 57 contrôleurs, authentification Sanctum avec portails séparés (\texttt{staff}/\texttt{client}), sessions uniques, RBAC maison (rôles + permissions \texttt{entité.action}) et périmètres de visibilité via \texttt{ScopeService}.

\begin{tabularx}{\textwidth}{|l|Y|}
\hline
\rowcolor{peklight}\textbf{Domaine} & \textbf{Ce qui existe} \\ \hline
Organisation & \texttt{organizations}, \texttt{countries}, \texttt{cities}, \texttt{agencies}, \texttt{departments}, chefs d'agence/département, \texttt{user\_assignments} \\ \hline
Catalogue & \texttt{categories}, \texttt{services} (+ promotions, historique prix, tiers séminaire), \texttt{products} (flag stock non exploité), \texttt{courses}, \texttt{training\_sessions}, \texttt{trainers}, \texttt{enrollments} \\ \hline
CRM & Clients = utilisateurs rôle \og client \fg{} (\texttt{client\_number}, catégorie, pays/ville, commercial référent), \texttt{prospects} + conversion, \texttt{commercials} (points, type commission) \\ \hline
Ventes & \texttt{invoices} + \texttt{invoice\_items} + \texttt{invoice\_payments} (cash/OM/MoMo), annulation, \texttt{orders}/\texttt{order\_lines}, vente rapide \\ \hline
Commissions & \texttt{commission\_payments} calculées à l'encaissement (taux snapshot par commercial), points fidélité \\ \hline
Abonnements & \texttt{subscription\_packs}, \texttt{subscriptions} (cycle de vie, renouvellement), \texttt{subscription\_notifications} (J-14/7/2/1) \\ \hline
Finance & \texttt{accounting\_categories}, \texttt{accounting\_transactions}, \texttt{daily\_balances}, bilans jour/période, statistiques (groupe, pays, agence, mensuel) \\ \hline
Système & \texttt{activity\_logs}, \texttt{login\_logs}, \texttt{settings}, uploads, exports CSV (maatwebsite/excel), Swagger L5 \\ \hline
\end{tabularx}

\subsection{Frontend (React 19, Vite, Tailwind v4, react-router 6)}

Kit UI maison (\texttt{components/ui/}), i18n FR/EN (~1530 clés), axios wrapper avec intercepteurs 401/403, contextes \texttt{AuthContext} et \texttt{ToastContext} uniquement.

Navigation actuelle : imbrication d'URLs \texttt{/} $\rightarrow$ \texttt{/countries/:id} $\rightarrow$ \texttt{/countries/:id/agencies/:id} $\rightarrow$ départements, chaque niveau exposant son contexte via \texttt{useOutletContext}. Un \texttt{AgencySwitcher} permet de changer d'agence depuis une agence, un \texttt{DepartmentSwitcher} entre départements d'une même agence.

\begin{tabularx}{\textwidth}{|l|Y|}
\hline
\rowcolor{peklight}\textbf{Module} & \textbf{État} \\ \hline
Dashboards & Groupe, pays, agence, département (pages implémentées) \\ \hline
Admin & Pays, agences, départements, utilisateurs, rôles/privileges, audit, settings, profil/2FA \\ \hline
Ventes/Factures & Liste, création, détail, encaissement, Quick Sale, impression ; comptabilité ; bilans ; abonnements \\ \hline
Academy & Cours, sessions, formateurs, apprenants (7 pages) — imbriqués sous l'\emph{agence} comme \og ligne Formations \fg{} \\ \hline
CRM & Clients (liste/détail) ; prospects : API prête mais \textbf{aucune page} \\ \hline
Store & \textbf{Absent} (aucune page, aucun appel API produits) \\ \hline
Studio & \textbf{Absent} totalement \\ \hline
Notifications & Dossier vide, aucune UI \\ \hline
\end{tabularx}

%=====================================================
\section{Analyse des écarts avec le cahier des charges}
%=====================================================

\begin{tabularx}{\textwidth}{|l|l|Y|}
\hline
\rowcolor{peklight}\textbf{Exigence (CDC)} & \textbf{État} & \textbf{Écart constaté} \\ \hline
Départements typés (4 branches) & Écart majeur & \texttt{departments} = simple libellé rattaché à une agence ; pas de colonne \texttt{type}. L'app a modélisé Academy comme une \og ligne \fg{} de l'agence (\texttt{agencies.type} agency/academy/mixed) au lieu d'un département. \\ \hline
Barre de contexte permanente & Manquant & Navigation par drill-down uniquement ; aucun sélecteur global Pays/Agence/Département/Période ; impossible de sauter directement vers une agence d'un autre pays. \\ \hline
Menus latéraux par branche & Partiel & Sidebar agence binaire (prestations/formations) ; les menus détaillés CDC (Store, Studio, Agency) n'existent pas. \\ \hline
Command Center + drill-down & Partiel & Stats existantes mais pas de vue Groupe consolidée KPI ni descente jusqu'à la transaction. \\ \hline
CRM complet & Partiel & Prospects sans UI, pas de pipeline visuel, pas d'activités/relances, pas d'entreprises. \\ \hline
Trésorerie multi-comptes & Partiel & Paiements liés aux factures mais aucun compte de trésorerie (caisse/MoMo/banque) ni mouvements rattachés. \\ \hline
Bilan journalier automatique & Partiel & \texttt{daily\_balances} saisis manuellement ; non généré depuis ventes/paiements/dépenses. \\ \hline
Dépenses avec workflow & Partiel & Dépenses = transactions comptables libres ; pas de cycle demande$\rightarrow$validation$\rightarrow$décaissement$\rightarrow$justificatif. \\ \hline
Commissions versionnées & Partiel & Taux fixé sur la fiche commercial ; aucune règle configurable/versionnée/paliers. \\ \hline
Agency : contrats/renouvellements & Partiel & Abonnements mensuels seulement ; pas d'objet contrat, publicité séparée du budget client. \\ \hline
Academy : présences/certificats/créances & Partiel & Présence booléenne sur inscription ; pas de feuille de présence par session ni certificats. \\ \hline
Store complet & Manquant & Produits (flag stock inutilisé) ; zéro stock, fournisseurs, achats, commandes store, livraisons, retours, inventaires. \\ \hline
Studio complet & Manquant & Rien : devis, projets, briefs, planning, production, révisions, livraisons. \\ \hline
Notifications centralisées & Manquant & Rappels d'abonnement uniquement. \\ \hline
Rapports/filtres communs & Partiel & Exports CSV ciblés ; pas de moteur de rapports filtrable période/pays/agence/département/commercial. \\ \hline
Migration Google Sheets & Manquant & Aucun outillage d'import. \\ \hline
\end{tabularx}

%=====================================================
\section{Recadrages d'architecture décidés}
%=====================================================

\subsection{Modèle organisationnel corrigé}

L'agence contient des \textbf{départements}, et chaque département porte un \textbf{type} parmi les quatre branches. Entrer dans une agence amène sur la page \textbf{Départements}, d'où l'on crée un département de l'un des quatre types. La notion d'\og agence academy/mixed \fg{} et le switcher \og Ligne Prestations / Ligne Formations \fg{} sont supprimés.

\begin{verbatim}
PEKEGNO GROUP
|-- CAMEROUN                        (country)
|   |-- DOUALA                      (agency)
|   |   |-- ACADEMY                 (department, type=academy)
|   |   |-- AGENCY                  (department, type=agency)
|   |   |-- STORE                   (department, type=store)
|   |   \-- STUDIO                  (department, type=studio)
|   \-- YAOUNDE                     (agency)
|       \-- ...
\-- COTE D'IVOIRE                   (country)
    \-- ABIDJAN                     (agency)
        \-- ...
\end{verbatim}

\subsection{Barre de contexte permanente}

Une barre supérieure unique, présente dans toutes les vues authentifiées :

\begin{center}
\texttt{PEKEGNO~|~Pays~$\blacktriangledown$~|~Agence~$\blacktriangledown$~|~Département~$\blacktriangledown$~|~Période~$\blacktriangledown$~|~Recherche~|~Notifications~|~Profil}
\end{center}

\begin{itemize}[nosep]
  \item Chaque sélecteur est alimenté par un arbre organisationnel chargé une seule fois (endpoint \texttt{/scope/context} enrichi) puis caché côté client.
  \item Sélection directe à n'importe quel niveau : depuis Cameroun $\rightarrow$ Yaoundé $\rightarrow$ Academy, on peut ouvrir \og Agence \fg{} et choisir directement une agence d'un autre pays — la barre recalcule le contexte et redirige proprement.
  \item Le contexte courant est persisté (localStorage + URL) et partagé via un \texttt{OrgContext} global ; les pages consomment le contexte au lieu de le recalculer depuis l'URL.
  \item La période (ce mois, ce trimestre, personnalisée) pilote tous les KPI affichés.
\end{itemize}

\subsection{Sidebars dynamiques par type de département}

Chaque type de département possède sa propre barre latérale, reprenant exactement les menus du cahier des charges :

\begin{description}[leftmargin=2.2cm, style=nextline]
  \item[type academy] Dashboard · Prospects · Apprenants · Inscriptions · Formations (Catalogue, Sessions, Planning, Formateurs) · Présences · Paiements · Créances · Certificats · Performances commerciales · Rapports
  \item[type agency] Dashboard · Prospects · Clients · Packages · Contrats · Prestations · Équipe client · Community Management · Publicité · Renouvellements · Paiements · Créances · Rapports clients · Performance
  \item[type store] Dashboard · Catalogue / Produits / Catégories · Stocks · Approvisionnements · Fournisseurs · Achats · Commandes · Ventes · Livraisons · Retours · Inventaires · Rapports
  \item[type studio] Dashboard · Prospects / Clients · Services · Devis · Projets · Briefs · Planning · Production · Révisions · Livraisons · Paiements · Rapports
\end{description}

Les routes migrent de \texttt{/countries/:c/agencies/:a/academy/*} vers \texttt{/departments/:departmentId/*} (redirections compatibles maintenues pendant la transition).

\subsection{Sémantique Studio (précisions métier)}

Un \textbf{projet} Studio est une production concrète commandée par une marque (ex. filmer une publicité). \textbf{Devis} : calcul du montant à verser pour le projet, éditable en lignes, soumis à validation puis conversion en projet après acompte. \textbf{Briefs} : réunions rattachées au projet ; chaque brief est daté, avec participants et \textbf{notes riches} saisies dans un éditeur de type Summernote (points de réunion, mise en forme, listes). \textbf{Planning} : plan de réalisation (jalons/tâches datées). \textbf{Production} : projets en cours de réalisation. \textbf{Révisions} : allers-retours avec le client avant validation. \textbf{Livraisons} : remise finale des livrables, solde, clôture.

%=====================================================
\section{Organisation de l'équipe}
%=====================================================

\subsection{Principe d'alternance}

Deux développeurs, \textbf{Blackstar} et \textbf{Le\_million}. Chaque sprint dure 2 semaines (10 jours ouvrés) et couvre un chantier. Les rôles \textbf{permutent à chaque sprint} afin que chacun progresse sur toute la stack :

\begin{tabularx}{\textwidth}{|c|l|c|c|}
\hline
\rowcolor{peklight}\textbf{Sprint} & \textbf{Semaines} & \textbf{Backend} & \textbf{Frontend} \\ \hline
S1 & S1--S2 & Blackstar & Le\_million \\ \hline
S2 & S3--S4 & Le\_million & Blackstar \\ \hline
S3 & S5--S6 & Blackstar & Le\_million \\ \hline
S4 & S7--S8 & Le\_million & Blackstar \\ \hline
S5 & S9--S10 & Blackstar & Le\_million \\ \hline
S6 & S11--S12 & Le\_million & Blackstar \\ \hline
S7 & S13--S14 & Blackstar & Le\_million \\ \hline
S8 & S15--S16 & Le\_million & Blackstar \\ \hline
S9 & S17--S18 & Blackstar & Le\_million \\ \hline
S10 & S19--S20 & Le\_million & Blackstar \\ \hline
\end{tabularx}

Capacité : $2 \times 10$ jours-homme par sprint, soit \textbf{200 j-h} sur 20 semaines.

\subsection{Definition of Done (toute tâche)}

\begin{itemize}[nosep]
  \item \textbf{BE} : migration testée (\texttt{migrate:fresh -{}-seed}), permissions \texttt{entité.action} câblées, policies/périmètres appliqués, endpoints documentés (Swagger régénéré), journalisation \texttt{activity\_logs} sur actions sensibles.
  \item \textbf{FE} : page branchée au menu latéral du bon scope, i18n FR/EN complète, états chargement/erreur/vide gérés, responsive mobile, permissions côté client alignées sur le backend.
  \item \textbf{QA} : parcours critique testé manuellement de bout en bout par le binôme (revue croisée), exports vérifiés.
\end{itemize}

%=====================================================
\section{Découpage des tâches par chantier}
%=====================================================

Durées en jours-homme (1 j $\approx$ 7 h). Resp. : BS = Blackstar, LM = Le\_million.

%-----------------------------------------------------
\subsection{C0 --- Recadrage organisationnel \& navigation contextuelle (Sprint 1)}
\addcontentsline{toc}{subsection}{C0 --- Recadrage organisationnel \& navigation contextuelle}

\begin{tabularx}{\textwidth}{|l|Y|c|c|c|}
\hline
\rowcolor{peklight}\textbf{ID} & \textbf{Tâche} & \textbf{Flux} & \textbf{Resp.} & \textbf{Durée} \\ \hline
T0.1 & Migration \texttt{departments.type} enum(\texttt{academy|agency|store|studio}) + backfill des données existantes, dépréciation de \texttt{agencies.type} & BE & BS & 3 j \\ \hline
T0.2 & Endpoint \texttt{/scope/context} enrichi : arbre pays$\rightarrow$agences$\rightarrow$départements (filtré par périmètre utilisateur), filtres par type & BE & BS & 3 j \\ \hline
T0.6 & Seeders de démonstration (agences avec départements des 4 types), Swagger, tests feature & BE & BS & 2 j \\ \hline
T0.3 & Store global de contexte (\texttt{OrgContext}) + \textbf{ContextBar} permanente : Pays/Agence/Département/Période + recherche + notifications (placeholder) + profil ; saut direct inter-pays/agences & FE & LM & 5 j \\ \hline
T0.4 & Suppression du switcher \og lignes \fg{} ; landing agence = page Départements ; modale de création de département avec les 4 types (icônes/couleurs) & FE & LM & 4 j \\ \hline
QA0 & Revue croisée parcours Groupe$\rightarrow$Pays$\rightarrow$Agence$\rightarrow$Département, redirections anciennes URLs, i18n & QA & Binôme & 3 j \\ \hline
\end{tabularx}

%-----------------------------------------------------
\subsection{C1 --- Trésorerie \& paiements (Sprint 2)}
\addcontentsline{toc}{subsection}{C1 --- Trésorerie \& paiements}

\begin{tabularx}{\textwidth}{|l|Y|c|c|c|}
\hline
\rowcolor{peklight}\textbf{ID} & \textbf{Tâche} & \textbf{Flux} & \textbf{Resp.} & \textbf{Durée} \\ \hline
T1.1 & Tables \texttt{treasury\_accounts} (cash / mobile money / banque, rattachées agence) + \texttt{treasury\_transactions} ; chaque paiement facture écrit un mouvement & BE & LM & 4 j \\ \hline
T1.2 & Bilan journalier \textbf{généré} : agrégation ventes/encaissements/dépenses/trésorerie, comparaison solde théorique vs réel, justification d'écart & BE & LM & 3 j \\ \hline
T1.7 & Tests, Swagger, seeders comptes (Caisses Douala/Yaoundé/Abidjan, OM, MTN MoMo, Afriland, CCA) & BE & LM & 1 j \\ \hline
T1.4 & Formulaire d'encaissement enrichi : sélection du compte de trésorerie destination, référence, affichage du compte sur l'historique paiements & FE & BS & 3 j \\ \hline
T1.5 & Page FINANCE/Trésorerie : soldes par compte, flux de la période, filtres compte/type/période, transferts internes (lecture v1) & FE & BS & 5 j \\ \hline
QA1 & Parcours : vente 100~000 FCFA en 3 tranches $\rightarrow$ 3 paiements $\rightarrow$ trésorerie $\rightarrow$ bilan cohérent & QA & Binôme & 4 j \\ \hline
\end{tabularx}

%-----------------------------------------------------
\subsection{C2 --- Dépenses \& C3 --- Commissions (Sprint 3)}
\addcontentsline{toc}{subsection}{C2 --- Dépenses \& C3 --- Commissions}

\begin{tabularx}{\textwidth}{|l|Y|c|c|c|}
\hline
\rowcolor{peklight}\textbf{ID} & \textbf{Tâche} & \textbf{Flux} & \textbf{Resp.} & \textbf{Durée} \\ \hline
T2.1 & Module dépenses : table dédiée, statuts brouillon$\rightarrow$soumis$\rightarrow$validé/rejeté$\rightarrow$payé/clos, catégories configurables, justificatifs (upload), notification validateur & BE & BS & 4 j \\ \hline
T2.2 & Moteur de commissions versionnées : \texttt{commission\_rules} (bénéficiaire, périmètre, formule \% / fixe / paliers, déclenchement, dates de validité), snapshot de règle sur chaque commission générée à l'encaissement & BE & BS & 4 j \\ \hline
QA2b & Tests règles multiples sur une même vente, reproductibilité des calculs & QA & BS & 2 j \\ \hline
T2.3 & Pages Dépenses : demande, file de validation (badge), décaissement lié au compte de trésorerie, pièces jointes, historique statuts & FE & LM & 5 j \\ \hline
T2.4 & Page COMMISSIONS : CRUD règles (formulaire paliers), liste commissions par période, workflow calculée$\rightarrow$à valider$\rightarrow$validée$\rightarrow$payée, traçabilité version de règle & FE & LM & 4 j \\ \hline
QA2 & Revue croisée workflow dépense + commission & QA & Binôme & 1 j \\ \hline
\end{tabularx}

%-----------------------------------------------------
\subsection{C4 --- CRM complet (Sprint 4)}
\addcontentsline{toc}{subsection}{C4 --- CRM complet}

\begin{tabularx}{\textwidth}{|l|Y|c|c|c|}
\hline
\rowcolor{peklight}\textbf{ID} & \textbf{Tâche} & \textbf{Flux} & \textbf{Resp.} & \textbf{Durée} \\ \hline
T4.1 & Opportunités/pipeline : stages contrôlés (nouveau, contacté, qualifié, proposition, négociation, gagné, perdu), valeur, date prévue & BE & LM & 3 j \\ \hline
T4.2 & Activités/relances : appels, RDV, notes, rappels planifiés + job quotidien de rappel & BE & LM & 3 j \\ \hline
T4.3 & Entprises (\texttt{companies}) + liaison contacts$\leftrightarrow$entreprises + recherche anti-doublons & BE & LM & 2 j \\ \hline
QA4b & Tests pipeline + rappels & QA & LM & 2 j \\ \hline
T4.4 & Kanban prospects/opportunités (drag \& drop entre stages, filtres commercial/agence/département) & FE & BS & 4 j \\ \hline
T4.5 & Fiche prospect complète (identité, acquisition, intérêt, responsable, pipeline, activités) + conversion client & FE & BS & 3 j \\ \hline
T4.6 & Timeline relationnelle client (achats, paiements, activités, relances) & FE & BS & 2 j \\ \hline
QA4 & Revue croisée CRM & QA & Binôme & 1 j \\ \hline
\end{tabularx}

%-----------------------------------------------------
\subsection{C5 --- Academy compléments \& C6 --- Contrats Agency (Sprint 5)}
\addcontentsline{toc}{subsection}{C5 --- Academy compléments \& C6 --- Contrats Agency}

\begin{tabularx}{\textwidth}{|l|Y|c|c|c|}
\hline
\rowcolor{peklight}\textbf{ID} & \textbf{Tâche} & \textbf{Flux} & \textbf{Resp.} & \textbf{Durée} \\ \hline
T5.1 & Feuille de présence par session (multi-apprenants, historique) remplaçant le booléen & BE & BS & 2 j \\ \hline
T5.2 & Certificats : génération (numéro, date, mention), statut délivré/retiré, PDF imprimable & BE & BS & 2 j \\ \hline
T6.1 & Contrats Agency : entité contrat (client, service/package, montant, cycles, échéance calculée), renouvellements avec alertes configurables, statuts actif / à renouveler / expiré / suspendu / résilié ; distinction honoraires vs budget publicitaire & BE & BS & 4 j \\ \hline
QA5b & Tests renouvellements + certificats & QA & BS & 2 j \\ \hline
T5.3 & UI présence session + certificats + fiche apprenant enrichie (échéancier N paiements, solde calculé, documents) & FE & LM & 5 j \\ \hline
T6.2 & Pages Contrats \& Renouvellements (alertes J-X, tâches de relance auto, historique) & FE & LM & 4 j \\ \hline
QA5 & Revue croisée Academy/Agency & QA & Binôme & 1 j \\ \hline
\end{tabularx}

%-----------------------------------------------------
\subsection{C7 --- Store / Matootoo (Sprints 6 et 7)}
\addcontentsline{toc}{subsection}{C7 --- Store / Matootoo}

\begin{tabularx}{\textwidth}{|l|Y|c|c|c|}
\hline
\rowcolor{peklight}\textbf{ID} & \textbf{Tâche} & \textbf{Flux} & \textbf{Resp.} & \textbf{Durée} \\ \hline
\multicolumn{5}{|l|}{\cellcolor{peklight}\textit{Sprint 6 --- S11--S12 (BE : LM / FE : BS)}} \\ \hline
T7.1 & Fournisseurs + achats/approvisionnements + réceptions (impact stock automatique) & BE & LM & 4 j \\ \hline
T7.2 & Stocks multi-agences : niveaux, \texttt{stock\_movements} typés (entrée, sortie, vente, retour, perte, ajustement, transfert inter-agence), seuils d'alerte & BE & LM & 4 j \\ \hline
QA7b & Tests mouvements + transferts & QA & LM & 2 j \\ \hline
T7.3 & UI Catalogue Store (produits/SKU/catégories) + page Stocks (niveaux, journal des mouvements, transfert guidé) & FE & BS & 5 j \\ \hline
T7.4 & UI Fournisseurs \& Achats (bon de commande, réception partielle) & FE & BS & 4 j \\ \hline
QA7.1 & Revue croisée & QA & Binôme & 1 j \\ \hline
\multicolumn{5}{|l|}{\cellcolor{peklight}\textit{Sprint 7 --- S13--S14 (BE : BS / FE : LM)}} \\ \hline
T7.5 & Commandes Store + livraisons + retours (impact stock et trésorerie, statuts contrôlés) & BE & BS & 4 j \\ \hline
T7.6 & Inventaires : sessions de comptage, écarts valorisés, ajustement traçable & BE & BS & 3 j \\ \hline
T7.7 & Rapports Store (rotation, marge, top produits) & BE & BS & 1 j \\ \hline
QA7c & Tests inventaire + retours & QA & BS & 2 j \\ \hline
T7.8 & Ventes Store : adaptation Quick Sale avec sortie de stock automatique & FE & LM & 3 j \\ \hline
T7.9 & UI Commandes / Livraisons / Retours & FE & LM & 4 j \\ \hline
T7.10 & UI Inventaires + Rapports Store & FE & LM & 2 j \\ \hline
QA7 & Revue croisée Store complet & QA & Binôme & 1 j \\ \hline
\end{tabularx}

%-----------------------------------------------------
\subsection{C8 --- Studio (Sprints 8 et 9)}
\addcontentsline{toc}{subsection}{C8 --- Studio}

\begin{tabularx}{\textwidth}{|l|Y|c|c|c|}
\hline
\rowcolor{peklight}\textbf{ID} & \textbf{Tâche} & \textbf{Flux} & \textbf{Resp.} & \textbf{Durée} \\ \hline
\multicolumn{5}{|l|}{\cellcolor{peklight}\textit{Sprint 8 --- S15--S16 (BE : LM / FE : BS)}} \\ \hline
T8.1 & Devis Studio : lignes, remises, TVA, statuts (brouillon, envoyé, accepté, refusé), conversion devis$\rightarrow$projet & BE & LM & 4 j \\ \hline
T8.2 & Projets (statuts pipeline), Briefs (réunions datées + notes riches Summernote côté FE), Planning (jalons/tâches datées) & BE & LM & 4 j \\ \hline
QA8b & Tests devis$\rightarrow$projet & QA & LM & 2 j \\ \hline
T8.3 & UI Devis (éditeur de lignes, aperçu imprimable) + liste Projets (pipeline visuel) & FE & BS & 5 j \\ \hline
T8.4 & UI Briefs (intégration éditeur Summernote) + Planning (vue calendrier/liste des jalons) & FE & BS & 4 j \\ \hline
QA8.1 & Revue croisée & QA & Binôme & 1 j \\ \hline
\multicolumn{5}{|l|}{\cellcolor{peklight}\textit{Sprint 9 --- S17--S18 (BE : BS / FE : LM)}} \\ \hline
T8.5 & Production/Révisions/Livrables : suivi d'avancement, cycles de révision, validation client, livraison, acompte/solde, clôture & BE & BS & 4 j \\ \hline
T9.1 & Centre de notifications unifié : table générale + règles déclencheurs (créance échue, renouvellement proche, stock sous seuil, dépense soumise, objectif) + destinataires configurables & BE & BS & 4 j \\ \hline
QA8c/QA9b & Tests notifications + clôture projet & QA & BS & 2 j \\ \hline
T8.6 & UI Production / Révisions / Livraisons + encaissements projet reliés au noyau de ventes & FE & LM & 5 j \\ \hline
T9.2 & Cloche de notifications (temps quasi-réel via polling), centre, préférences & FE & LM & 4 j \\ \hline
QA8/QA9 & Revue croisée & QA & Binôme & 1 j \\ \hline
\end{tabularx}

%-----------------------------------------------------
\subsection{C9 --- Command Center, rapports \& préparation migration (Sprint 10)}
\addcontentsline{toc}{subsection}{C9 --- Command Center, rapports \& préparation migration}

\begin{tabularx}{\textwidth}{|l|Y|c|c|c|}
\hline
\rowcolor{peklight}\textbf{ID} & \textbf{Tâche} & \textbf{Flux} & \textbf{Resp.} & \textbf{Durée} \\ \hline
T10.1 & API KPI consolidée : CA, encaissements, créances, dépenses, trésorerie, performance ; drill-down Groupe$\rightarrow$Pays$\rightarrow$Agence$\rightarrow$Département$\rightarrow$transaction ; comparaison période précédente & BE & LM & 4 j \\ \hline
T10.2 & Exports Excel/PDF/CSV étendus + squelette importeur CSV Google Sheets (mapping configurable, journal d'anomalies, import test) & BE & LM & 4 j \\ \hline
QA10b & Cohérence KPI vs données sources & QA & LM & 2 j \\ \hline
T10.3 & Command Center : vue Groupe consolidée, cartes KPI cliquables avec descente jusqu'à la transaction, alertes & FE & BS & 5 j \\ \hline
T10.4 & RAPPORTS transverses : filtres communs période/pays/agence/département/commercial, exports & FE & BS & 3 j \\ \hline
T10.5 & Recette générale du périmètre, corrections, documentation & QA & Binôme & 2 j \\ \hline
\end{tabularx}

%-----------------------------------------------------
\subsection{Synthèse des charges}

\begin{tabularx}{\textwidth}{|l|Y|c|c|}
\hline
\rowcolor{peklight}\textbf{Chantier} & \textbf{Contenu} & \textbf{Charge} & \textbf{Calendaire} \\ \hline
C0 & Organisation, ContextBar, départements typés & 20 j-h & S1--S2 \\ \hline
C1 & Trésorerie, bilan journalier auto & 20 j-h & S3--S4 \\ \hline
C2/C3 & Dépenses workflow, commissions versionnées & 20 j-h & S5--S6 \\ \hline
C4 & CRM complet (pipeline, activités, entreprises) & 20 j-h & S7--S8 \\ \hline
C5/C6 & Academy compléments, contrats \& renouvellements & 20 j-h & S9--S10 \\ \hline
C7 & Store complet (stocks, achats, ventes, inventaires) & 40 j-h & S11--S14 \\ \hline
C8/C10 & Studio complet + notifications & 40 j-h & S15--S18 \\ \hline
C9 & Command Center, rapports, outillage migration & 20 j-h & S19--S20 \\ \hline
\rowcolor{peklight}\textbf{Total} & & \textbf{200 j-h} & \textbf{20 semaines} \\ \hline
\end{tabularx}

%=====================================================
\section{Calendrier prévisionnel}
%=====================================================

\begin{figure}[h!]
\centering
\scriptsize
\begin{tabular}{|l|*{20}{W}|}
\hline
\rowcolor{peklight}\textbf{Chantier / Semaine} & 1 & 2 & 3 & 4 & 5 & 6 & 7 & 8 & 9 & 10 & 11 & 12 & 13 & 14 & 15 & 16 & 17 & 18 & 19 & 20 \\ \hline
C0 Navigation \& contexte & \cellcolor{pekblue} & \cellcolor{pekblue} & & & & & & & & & & & & & & & & & & \\ \hline
C1 Trésorerie \& paiements & & & \cellcolor{pekblue} & \cellcolor{pekblue} & & & & & & & & & & & & & & & & \\ \hline
C2 Dépenses / C3 Commissions & & & & & \cellcolor{pekgreen} & \cellcolor{pekgreen} & & & & & & & & & & & & & & \\ \hline
C4 CRM complet & & & & & & & \cellcolor{pekgreen} & \cellcolor{pekgreen} & & & & & & & & & & & & \\ \hline
C5 Academy / C6 Contrats & & & & & & & & & \cellcolor{pekgreen} & \cellcolor{pekgreen} & & & & & & & & & & \\ \hline
C7 Store & & & & & & & & & & & \cellcolor{pekamber} & \cellcolor{pekamber} & \cellcolor{pekamber} & \cellcolor{pekamber} & & & & & & \\ \hline
C8 Studio + Notifications & & & & & & & & & & & & & & & \cellcolor{pekamber} & \cellcolor{pekamber} & \cellcolor{pekamber} & \cellcolor{pekamber} & & \\ \hline
C9 Command Center \& rapports & & & & & & & & & & & & & & & & & & & \cellcolor{pekblue} & \cellcolor{pekblue} \\ \hline
\rowcolor{peklight}\textbf{Jalon} & \multicolumn{2}{c|}{J1} & \multicolumn{2}{c|}{J2} & \multicolumn{2}{c|}{J3a} & \multicolumn{2}{c|}{J3b} & \multicolumn{2}{c|}{J3c} & \multicolumn{4}{c|}{J4} & \multicolumn{4}{c|}{J5} & \multicolumn{2}{c|}{J6} \\ \hline
\end{tabular}
\caption{Planning calendaire (1 cellule = 1 semaine). J1 : navigation recadrée · J2 : noyau financier · J3 : métiers prioritaires complets (fin S10) · J4 : Store livré · J5 : Studio + notifications · J6 : recette finale.}
\end{figure}

%=====================================================
\section{Jalons de validation}
%=====================================================

\begin{tabularx}{\textwidth}{|l|Y|c|}
\hline
\rowcolor{peklight}\textbf{Jalon} & \textbf{Contenu validé avec Pekegno} & \textbf{Fin} \\ \hline
J1 & Arbre pays$\rightarrow$agence$\rightarrow$départements typés opérationnel, ContextBar, menus latéraux par branche & S2 \\ \hline
J2 & Trésorerie multi-comptes, bilan journalier auto, dépenses workflow, commissions versionnées & S6 \\ \hline
J3 & CRM complet, Academy complété, contrats/renouvellements Agency & S10 \\ \hline
J4 & Store complet (catalogue, stocks, achats, ventes, livraisons, retours, inventaires) & S14 \\ \hline
J5 & Studio complet + centre de notifications & S18 \\ \hline
J6 & Command Center drill-down, rapports transverses, recette, outillage migration prêt & S20 \\ \hline
\end{tabularx}

%=====================================================
\section{Points de décision bloquants (à arbitrer avec Pekegno)}
%=====================================================

Ces arbitrages conditionnent les tâches indiquées et doivent être tranchés \textbf{avant} le sprint concerné :

\begin{enumerate}[nosep]
  \item Définition comptable exacte du chiffre d'affaires par activité (avant S19, T10.1).
  \item Moment de déclenchement des commissions et formules réelles de prime : \% / fixe / paliers (avant S5, T2.2).
  \item Workflow et seuils de validation des dépenses (avant S5, T2.1).
  \item Règles de renouvellement Agency et délais d'alerte (avant S9, T6.1).
  \item Liste définitive des comptes de trésorerie Cash / Mobile Money / banques (avant S3, T1.1).
  \item Matrice définitive rôles/périmètres incluant les départements typés (avant S1, T0.2).
  \item Périmètre et qualité des données Sheets 2022--2026 à migrer (avant S19, T10.2).
\end{enumerate}

%=====================================================
\section{Risques et mitigations}
%=====================================================

\begin{tabularx}{\textwidth}{|Y|Y|}
\hline
\rowcolor{peklight}\textbf{Risque} & \textbf{Mitigation} \\ \hline
KPI non figés par la direction financière & Afficher dès J2 des définitions provisoires paramétrables dans \texttt{settings} ; figer au plus tard fin S18. \\ \hline
Formules de commissions complexes (paliers) & Moteur de règles générique + snapshot versionné : une règle nouvelle n'altère jamais les commissions passées. \\ \hline
Refonte navigation perturbe les utilisateurs existants & Redirections permanentes des anciennes URLs pendant tout le chantier C0 + revue avec Pekegno au jalon J1. \\ \hline
Alternance BE/FE ralentit un développeur sur un domaine & DoD stricte, revue croisée systématique en QA, documentation Swagger et types TS synchronisés à chaque tâche BE. \\ \hline
Données Sheets de mauvaise qualité & Import test obligatoire avec rapprochement des totaux historiques ; anomalies tracées plutôt que converties silencieusement. \\ \hline
Dérive du périmètre Store/Studio & Menus latéraux figés par le cahier des charges ; toute addition passe par une décision documentée. \\ \hline
\end{tabularx}

%=====================================================
\section{Critères de fin de projet}
%=====================================================

\begin{itemize}[nosep]
  \item Une donnée métier est saisie une seule fois ; soldes, créances, trésorerie et commissions en découlent automatiquement.
  \item Vente $\neq$ paiement respectée partout (N paiements par vente, aucun champ \og tranche 1/tranche 2 \fg{}).
  \item Depuis la barre de contexte, tout utilisateur autorisé atteint n'importe quel pays/agence/département en une interaction.
  \item Chaque département affiche exactement les menus de sa branche.
  \item Le bilan journalier se génère sans ressaisie ; les écarts se justifient dans l'outil.
  \item Les opérations sensibles (paiements, dépenses, commissions, permissions) sont auditables.
  \item L'importeur Sheets exécute un import test dont les totaux sont rapprochés des sources.
\end{itemize}

\end{document}
